\chapter{Finite resolution, analytic error budgets and laboratory limits}
\label{ch:finite-resolution}

\reviewstatus{Release dated 13 November 2026, America/Los\_Angeles.
This chapter supplies written truncation bounds for three declared analytic
families and connects them to the existing laboratory kernels. Independent
80-decimal-digit reference calculations and focused numerical tests check
selected instances. The algebraic square-perturbation component is separately
identified in the Lean companion; the trigonometric estimates, continuum
sampling assumptions and floating-point evaluations are not promoted to
Lean or interval certificates. A finite-box diffusion gap, a finite-field
example and a sampled exact flow have different meanings. No continuum
regularity, mass-gap, arbitrary-field convergence or physical quantum
advantage result follows from these diagnostics.}

\section{Which limit is being tested?}

A refinement study must hold a physical problem fixed while changing a
declared approximation. In the flow and heat panels the periodic box has
side length $L=2\pi$ and the mesh width is $h=L/n$. The gauge panel
instead uses one centered cell of side $h=1/n$ inside a unit-square patch.
Its central curvature and quadratic variation are held fixed as the cell
shrinks. These are different families; the old single-plaquette flux
slider has no mesh width of its own.

\begin{longtable}{p{.22\textwidth}p{.67\textwidth}}
\toprule Quantity & Meaning and limitation\\\midrule
\endhead
Analytic truncation & A proved bound in real arithmetic, conditional on the
stated field, smoothness, grid and step restrictions.\\
Observed error & A difference between evaluated formulas or kernels.
It can reflect cancellation and library arithmetic as well as truncation.\\
Rounding indicator & A heuristic scale involving machine precision and
conditioning. It is not a directed-rounding enclosure.\\
Finite-volume effect & Dependence on the physical box size $L$.
Refining $n$ at fixed $L$ does not test $L\to\infty$.\\
Model discrepancy & A difference between a chosen family and the intended
physical system, boundary conditions, measured parameters or initial data.
Mesh refinement alone does not estimate it.\\
\bottomrule
\end{longtable}

For two errors $E_h,E_{h/2}>0$, the observed rate is
$\log(E_h/E_{h/2})/\log2$. It is an observation on that pair of grids.
The implementation withholds it when either error is at its heuristic
rounding scale or a reference has underflowed. It does not force an expected
rate by selecting favorable grids. In particular, errors from two
approximations may cancel, making their sum appear unusually accurate.

The new functions in \path{extension/site/finite-resolution-core.js}
reuse the existing \texttt{taylorGreen}, \texttt{heatStep} and
\texttt{gaugePlaquette} implementations. They introduce no new time-marching
Navier--Stokes solver. The adjacent panels have explicitly separate
diagnostic controls, so animation or a finite certificate import cannot
silently change the meaning of a refinement table.

\section{Taylor--Green derivatives and a nonlinear residual}

Write $a=e^{-2\nu t}$, with $\nu,t\ge0$, and use the existing solution
\[
 u=a\sin x\cos y,\qquad v=-a\cos x\sin y,\qquad
 p=\frac{a^2}{4}(\cos2x+\cos2y).
\]
All quantities are sampled on the $n\times n$ periodic grid. We require
$n\ge5$: the velocity frequency one and pressure frequency two are then
strictly below Nyquist. Let $D_{x,h},D_{y,h}$ be centered first differences,
$\Delta_h$ the five-point Laplacian, and
$s(z)=\sin z/z$, with $s(0)=1$. Exact Fourier calculation gives
\[
 D_{x,h}e^{ikx}=ik\,s(kh)e^{ikx},\qquad
 -\Delta_he^{i(k_xx+k_yy)}
 =\frac4{h^2}\left(\sin^2\frac{k_xh}{2}
                   +\sin^2\frac{k_yh}{2}\right)e^{i(k_xx+k_yy)}.
\]
Here and below the maximum norm of a vector takes the largest absolute
component over the sampled grid.

\begin{proposition}[Bounds for the sampled manufactured flow]
\label{prop:fr-taylor-green}
Put $\lambda_h=8\sin^2(h/2)/h^2$ and
$\omega=2a\sin x\sin y$. In real arithmetic,
\begin{align*}
 \|D_{x,h}v-D_{y,h}u-\omega\|_\infty
     &\le \frac{ah^2}{3},\\
 \|\Delta_h(u,v)+2(u,v)\|_\infty
     &\le \frac{ah^2}{6},\\
 D_{x,h}u+D_{y,h}v&=0.
\end{align*}
With the time derivative evaluated analytically, define
\[
 R_h=\partial_t(u,v)
   +uD_{x,h}(u,v)+vD_{y,h}(u,v)
   +(D_{x,h}p,D_{y,h}p)-\nu\Delta_h(u,v).
\]
Then
\[
 \|R_h\|_\infty\le \frac{\nu ah^2}{6}+\frac{a^2h^2}{4}.
\]
\end{proposition}
\begin{proof}
Taylor's integral remainder gives
$|1-s(z)|\le z^2/6$ and
$0\le1-4\sin^2(z/2)/z^2\le z^2/12$ for $z\ne0$.
Thus the computed curl is $s(h)\omega$,
$0\le2-\lambda_h\le h^2/6$, and the first two bounds follow.
The two centered divergence terms are
$as(h)\cos x\cos y$ and its negative.

The advective and pressure derivatives carry different Fourier
multipliers. Direct substitution gives
\[
 \begin{split}
 R_{h,x}&=\nu(\lambda_h-2)u+
       \frac{a^2}{2}\sin2x\,\{s(h)-s(2h)\},\\
 R_{h,y}&=\nu(\lambda_h-2)v+
       \frac{a^2}{2}\sin2y\,\{s(h)-s(2h)\}.
 \end{split}
\]
Since $0<h\le2\pi/5<\pi$, we have $0\le s(h)\le1$ and
\[
 0\le s(h)-s(2h)=s(h)(1-\cos h)\le h^2/2.
\]
Taking maximum component norms proves the residual estimate.
\end{proof}

This calculation tests a nonlinear product and a pressure cancellation,
not just a linear derivative. Nevertheless it measures consistency on
one known solution. A small residual does not by itself bound the error
of an arbitrary numerical trajectory. Such a conclusion also needs
stability, compatible initial data and control of the relevant solution
class.

\begin{proposition}[Exact energy quadrature for this Fourier family]
\label{prop:fr-energy-quadrature}
For every grid admitted above, the real-arithmetic average
\[
 E_n=\frac1{n^2}\sum_{j,\ell=0}^{n-1}
       \frac{u(jh,\ell h)^2+v(jh,\ell h)^2}{2}
\]
equals the continuum average $a^2/4$.
\end{proposition}
\begin{proof}
The energy density is
$a^2(1-\cos2x\cos2y)/4$. The geometric sum
$\sum_{j=0}^{n-1}e^{4\pi i j/n}$ vanishes for $n\ge5$,
so the product cosine term has zero grid average. Its integral has zero
average as well.
\end{proof}

The exactness is a finite Fourier identity. More general periodic
quadrature is governed by regularity and aliasing; the Fourier explanation
in \citet[Section~3]{frTrefethenWeideman2014} supplies broader context.
We do not infer exact quadrature for an arbitrary smooth velocity field.
The implementation uses compensated summation and still records the
floating-point energy discrepancy.

\section{Heat at fixed physical time}

Let $k=(k_x,k_y)\in\Z^2$, with $|k_x|,|k_y|<n/2$, and initially take
$k\ne0$. On the same torus use the normalized grid probability
\[
 p^0_{j\ell}=\frac{1+\frac12\cos(k_xjh+k_y\ell h)}{n^2}.
\]
The zero mode is treated separately as the uniform probability. Let
$T=m\Delta t$, where $m$ is a positive integer when $T>0$, and require
\[
 r=\frac{\nu\Delta t}{h^2}\le\frac14.
\]
The existing heat step is the convex combination
\[
 p^{s+1}_{j\ell}=(1-4r)p^s_{j\ell}
    +r(p^s_{j+1,\ell}+p^s_{j-1,\ell}
        +p^s_{j,\ell+1}+p^s_{j,\ell-1}).
\]
Define
\[
 \lambda=|k|^2,\qquad
 \lambda_h=\frac4{h^2}
   \left(\sin^2\frac{k_xh}{2}+\sin^2\frac{k_yh}{2}\right),
 \quad x=\nu\Delta t\lambda_h,\quad q=1-x.
\]
The discrete, spatially discrete continuous-time, and continuum mode
amplitudes are, respectively,
$A_d=q^m$, $A_s=e^{-\nu T\lambda_h}$ and $A_c=e^{-\nu T\lambda}$.
For $T=0$ all are one and the step count is zero.

\begin{proposition}[Separated spatial and time budgets]
\label{prop:fr-heat-mode}
Under these assumptions,
\begin{align*}
 |A_s-A_c|&\le
    \min\left(1,\frac{\nu T h^2}{12}(k_x^4+k_y^4)\right)=:B_s,\\
 |A_d-A_s|&\le\min\left(2,\frac{mx^2}{2}\right)=:B_t,\\
 |A_d-A_c|&\le\min(2,B_s+B_t).
\end{align*}
The formulas and bounds remain valid when $q<0$.
\end{proposition}
\begin{proof}
The trigonometric remainder used above gives
$0\le\lambda-\lambda_h\le h^2(k_x^4+k_y^4)/12$.
The map $z\mapsto e^{-\nu Tz}$ has derivative of magnitude at most
$\nu T$ for $z\ge0$, and its values lie in $[0,1]$. This proves the
first bound.

Since $0\le x\le8r\le2$, we have $|q|\le1$. For $b=e^{-x}$,
\[
 |q^m-b^m|
 \le |q-b|\sum_{j=0}^{m-1}|q|^j|b|^{m-1-j}
 \le m|q-b|.
\]
The integral remainder
$e^{-x}-(1-x)=\int_0^x(x-s)e^{-s}\,ds$
lies between zero and $x^2/2$. The amplitudes also have absolute value at
most one, proving the second bound with its cap of two. The triangle
inequality proves the third. No step uses $q\ge0$.
\end{proof}

An admissible negative multiplier changes the sign of a Fourier component;
it does not make this nonnegative initial probability become negative.
The channel remains a convex combination of translations. At the
positivity boundary an even-grid checkerboard can have multiplier $-1$,
so positivity alone does not imply strict convergence to equilibrium.
The implemented diagnostic tests an admissible negative-mode case and
separately rejects an unstable CFL ratio.

The probability discrepancy at a grid point is at most
$|A_d-A_c|/(2n^2)$ and its discrete $\ell^1$ norm is at most
$|A_d-A_c|/2$. These conversion factors matter when an amplitude bound
is used in an observable or control norm.

The default refinement chooses
$m=\max(1,\lceil5\nu T/h^2\rceil)$ for $T>0$ and $m=0$ at $T=0$.
It keeps $r\le1/5$ and holds $T$ fixed. For $\nu T>0$,
\[
 \frac{mx^2}{2}
 =\frac{\nu^2T^2\lambda_h^2}{2m}
 \le\frac{\nu T h^2\lambda^2}{10},
\]
so the sum of the two analytic budgets is $O(h^2)$ for a fixed mode.
Holding the original time step fixed at all grid sizes would eventually
violate the positivity restriction.

The two signed errors can oppose one another: the spatial Laplacian
underestimates the decay rate, whereas forward Euler can overestimate
decay for a low-frequency stable mode. With $\nu=0.08$, $T=1$ and
$k=(1,0)$, the recorded total amplitude errors at
$n=8,16,32,64$ are approximately
$9.12\times10^{-4}$, $3.26\times10^{-5}$,
$3.12\times10^{-5}$ and $1.10\times10^{-5}$.
The intervening observed rates are about $4.805$, $0.067$ and $1.501$.
They are retained alongside the separate spatial and time errors.
They neither refute the proved upper bound nor establish a higher
convergence order.

The point and peak fixtures in the original heat view use grid-index
widths. Varying $n$ changes their physical width. The smooth Fourier
family above therefore provides a declared consistency test; no
fixed-data continuum limit is inferred for those other fixtures.

\subsection{A finite box has its own diffusion gap}

For a square periodic box of arbitrary side $L$, the continuum heat
generator's smallest positive eigenvalue for $\nu>0$ is
$g_L=\nu(2\pi/L)^2$. The corresponding five-point generator value is
\[
 g_{L,n}=\frac{4\nu}{(L/n)^2}\sin^2\frac{\pi}{n},\qquad
 0\le1-\frac{g_{L,n}}{g_L}
     =1-\left(\frac{\sin(\pi/n)}{\pi/n}\right)^2
     \le\frac{\pi^2}{3n^2}.
\]
This follows from the discrete Fourier eigenvalues and the same cosine
remainder. The diagnostic reports both values for $L=2\pi$ and uses
$0\le g_L-g_{L,n}\le\nu h^2/12$; at $\nu=0$ both gaps are zero.
These are generator eigenvalues, not an absolute mixing gap for every
discrete-time channel.

At fixed $L$, refinement approaches a positive continuum box gap.
Increasing $L$ instead sends $g_L$ to zero. Thus a finite-volume diffusion
gap does not establish a volume-uniform lower bound, and it has no
automatic interpretation as a quantum Yang--Mills mass gap.

\section{A local U(1) plaquette with declared curvature}

Wilson's lattice construction motivates retaining a gauge-invariant loop
quantity while examining the small-cell limit. The classical lattice
approximation and the quantized continuum limit already require separate
arguments in that construction.\cite[Section~III.A]{frWilson1974}
The present example is the explicitly Abelian local calculation below.

On $C_h=[-h/2,h/2]^2$, set
\[
 B(x,y)=B_0+\kappa(x^2+y^2),\qquad
 \overline B_h=B_0+\frac{\kappa h^2}{6},\qquad
 \Phi_h=\int_{C_h}B\,dx\,dy=h^2\overline B_h.
\]
There is no edge-integration approximation. Reusing the existing four-link
product gives $W=e^{i\Phi_h}$ in real arithmetic. The normalized local
action density is
\[
 Q_h=\frac{2(1-\Re W)}{h^4}
     =\frac{4\sin^2(\Phi_h/2)}{h^4}.
\]
The target is the central value $B_0^2$, not a spatially averaged field
energy over the entire unit square.

\begin{proposition}[Local cell averaging and small-angle budget]
\label{prop:fr-local-wilson-bound}
Suppose an exact cell mean obeys
$|\overline B_h-B_0|\le d$. Then
\[
 |Q_h-B_0^2|
 \le d(2|B_0|+d)+\frac{h^4}{12}(|B_0|+d)^4.
\]
For the displayed quadratic field one may take
$d=|\kappa|h^2/6$. More generally, for a twice continuously
differentiable scalar field on $C_h$ with
$|B_{xx}|\le M_{xx}$ and $|B_{yy}|\le M_{yy}$, one may take
$d=h^2(M_{xx}+M_{yy})/24$.
\end{proposition}
\begin{proof}
For real $z$, Taylor's remainder and $1-\cos z\le z^2/2$ give
$0\le z^2-2(1-\cos z)\le z^4/12$. Hence
\[
 |Q_h-\overline B_h^2|
 \le h^4|\overline B_h|^4/12
 \le h^4(|B_0|+d)^4/12.
\]
Writing $e=\overline B_h-B_0$, the factorization
$(B_0+e)^2-B_0^2=e(2B_0+e)$ gives
$|\overline B_h^2-B_0^2|\le d(2|B_0|+d)$.
The triangle inequality proves the asserted budget.

For the quadratic field, integrate $x^2+y^2$ over the square.
For the smooth-field version, the one-dimensional midpoint remainder is
bounded in average by $Mh^2/24$ because the linear term integrates to
zero and the quadratic remainder is at most $Mx^2/2$. Apply this
first in $y$ at each $x$ and then in $x$ at $y=0$.
The two resulting bounds add; no mixed derivative bound is needed.
\end{proof}

The freshly kernel-checked formal component
\path{LaboratoryContracts.square_perturbation_bound}
encodes exactly the factorization-based scalar estimate from
$|e|\le d$. Its role is a component of
Proposition~\ref{prop:fr-local-wilson-bound}; it does not prove the
smoothness assumptions, cell integration, cosine remainder or numerical
evaluation. If a separate flux approximation had a certified error
$|\widehat\Phi_h-\Phi_h|\le\eta_\Phi$, the same argument could use
$d+\eta_\Phi/h^2$ for its inferred cell mean. That additional error is
zero in the implemented exact-flux family and is not estimated for a
general measured field.

For $\kappa=0$ the bound is $O(h^4)$; with a varying field it is
generally $O(h^2)$. A special vanishing central field can remove the
leading term, so these generic rates are not lower bounds on error.
Likewise, $W$ only records flux modulo $2\pi$. The examples with
fluxes $8$ and $8+2\pi$ have the same Wilson quantity and different
unwrapped $B_0^2$ at $h=1$. The bound never resolves that ambiguity.
The UI flags $|\Phi_h|\ge\pi$ and retains this negative control.

\subsection{Cancellation is visible before it is repaired}

The raw implementation forms a product of four complex link phases and
subtracts its real part from one. The equivalent sine-squared expression
avoids this subtraction for a small flux. At $B_0=1$, $\kappa=0$ and
$n=65536$, the recorded binary64 raw density is zero, while the stable
value rounds to one. The real analytic deficit is positive and about
$4.52\times10^{-21}$; the 80-digit reference retains it. Therefore
the stable zero error at this scale receives an unresolved rate rather
than a claim of perfect convergence.

The raw error can be amplified on the scale
$\epsilon_{\rm mach}/h^4$. The code exposes a deliberately labeled
heuristic multiple of that scale, alongside a heuristic for the stable
formula. These are diagnostic estimates, not verified outward-rounded
enclosures. The ECMAScript specification makes the relevant ordinary
sine, cosine and exponential evaluations implementation-approximated;
it does not turn these formulas into rigorous intervals.\cite{frECMAScriptLiving}

\section{Exact topology, finite arithmetic and nonlinear stability}

An exact invariant and a sampling error should not share an unlabeled
error bar. The existing finite cell complex satisfies its boundary
identity algebraically; the Hodge parameter family has an exact kernel
change at zero; and the integer matrix $J^k$ has determinant $14^k$.
Rasterization or display quadrature does not change those statements.
A numerical rank inferred from a tolerance, however, may change when
the tolerance or conditioning changes. The exact laboratory bridge
therefore uses rational or integer inputs for its finite contracts.

The elliptic-curve example over a fixed finite field is an exact finite
model. It is not a grid approximation to a real fluid or gauge field.
Similarly, refining the drawn torus does not change its algebraic
covering degree. A continuum relation would require a specified sequence
of objects, comparison maps and an explicit convergence theorem.

There is a useful conditional connection from a residual to a nonlinear
trajectory error. Suppose $\dot y=F(y)$, the surrogate $z$ is absolutely
continuous, and both paths stay in a region on which $F$ has Lipschitz
constant $K\ge0$ in a specified norm. If
$r=\dot z-F(z)$ is integrable, the integral equations and Gronwall's
inequality give
\[
 \|z(t)-y(t)\|
 \le e^{Kt}\|z(0)-y(0)\|
      +\int_0^t e^{K(t-s)}\|r(s)\|\,ds.
\]
Indeed the triangle inequality gives the corresponding integral
inequality for the error; multiplying the usual comparison function
by $e^{-Kt}$ and integrating yields the displayed estimate.
It is a conditional finite-dimensional stability statement.
The Taylor--Green grid residual alone supplies neither the common
invariant region nor a continuum PDE Lipschitz estimate.

For a contractive feedback realization, the appropriate bookkeeping is
different. If a reference map $T$ has contraction factor $q<1$ in the
chosen weighted norm and an implemented map obeys a \emph{uniform}
one-step error bound $\|\widetilde T(x)-T(x)\|\le\eta$, then
\[
 e_{m+1}\le qe_m+\eta,\qquad
 e_m\le q^m e_0+\frac{1-q^m}{1-q}\eta.
\]
The second formula follows by iterating the first geometric sum.
The required $\eta$ may include spatial, time, arithmetic, encoding
and model errors, but only after they are expressed in the same norm
and with compatible deterministic or probabilistic hypotheses.
An error verified on one Fourier mode cannot be inserted as a uniform
bound for all states. A heuristic rounding scale cannot be inserted
as a certified summand. These conditions explain how the finite
weighted-gain results in Chapter~\ref{ch:quantum-tropical-realizations}
can consume a valid numerical budget without absorbing unsupported
continuum or physical assumptions.

\section{Implementation evidence and reproduction}

The independent reference generator uses 80 decimal digits and exact
spectral expressions, including the nonlinear residual symbol. It
retains both binary64 comparison values and longer decimal strings.
The focused numerical verifier checks 24 reference cases,
24 Taylor--Green grid cases, 36 heat-budget cases and 23 invalid-input
controls. It also compares repeated calls to the existing heat kernel,
tests the constant and varying-curvature gauge families, and preserves
phase aliases, raw cancellation, underflow and zero-error cases.
The largest recorded reference discrepancies are approximately
$4.98\times10^{-15}$ for the sampled flow,
$1.11\times10^{-15}$ for heat and
$2.57\times10^{-16}$ for the gauge scalar formulas.
These figures describe the retained run, not uniform arithmetic bounds.

The isolated UI verifier checks eight behavioral groups using an
in-memory DOM fixture: labeled controls and table captions, invalid-input
clearing, recovery, event updates and visible numerical limitations.
It does not claim a rendered-browser, keyboard or screen-reader audit.
Exact run dates in the machine-readable records are labeled UTC.

\begin{lstlisting}
python3 extension/code/build_finite_resolution_reference.py
node extension/site/verify-finite-resolution.cjs
node extension/site/verify-finite-resolution-ui.cjs
\end{lstlisting}

The numerical and UI receipts are
\path{extension/verification/finite-resolution-node.json} and
\path{extension/verification/finite-resolution-ui.json}.
The source ledger and exact reading scope are in
\path{extension/review/finite-resolution-20261113.md}.
These records keep the analytic contract, implementation comparison,
formal component and remaining physical assumptions separately inspectable.
